% Prose rendering of PrimeTensor/MulRat.lean.
% Fragment: no \documentclass, no preamble, no document environment.

\section{Positive rational magnitudes as barcodes}
\label{MulRat:sec}

\leanfile{PrimeTensor/MulRat.lean}

\subsection*{What this file establishes}

\texttt{PrimeTensor/Ratio.lean} introduced $\name{PrimeRatio}$, an oriented
pair of prime multisets, and defined the cross-product relation
$\name{Same}$ on it --- but proved nothing whatever about that relation.
This file supplies everything that was missing and then takes the quotient.

The result, $\name{MulRat}$, is a commutative group: multiplication is
associative and commutative with unit $1$, and every element has an inverse.
All six laws are equalities of elements, not statements modulo a relation.

One step in the file does real work.  Reflexivity and symmetry of
$\name{Same}$ are immediate, but transitivity requires cancelling a common
factor, and cancellation on $\mathbb{N}$ is valid only against a nonzero
factor.  The factor available is the evaluation of a prime multiset, and it is
nonzero because no bag evaluates below $1$.  This is the point at which the
project's refusal of a zeroth stage, decided in
\texttt{PrimeTensor/Depth.lean} and carried into
\texttt{PrimeTensor/Prime.lean}, is actually spent.

Finally, and as the module header states, no value of type $\mathbb{Q}$ or
$\mathbb{R}$ occurs anywhere in the file.  Magnitudes are barcodes throughout;
$\mathbb{N}$ enters only through $\name{eval}$, and only to compare.

\subsection{Positivity}

\begin{lemma}[\lean{PrimeRatio.multiset\_eval\_pos}]\label{MulRat:lem:pos}
$0 < \name{eval}(q)$ for every $q : \name{PrimeMultiset}$.
\end{lemma}

\begin{leancode}
private theorem multiset_eval_pos (q : PrimeMultiset) : 0 < q.eval := by
  refine Quotient.inductionOn q ?_
  intro b
  exact lt_of_lt_of_le Nat.zero_lt_one (PrimeBag.one_le_eval b)
\end{leancode}

\begin{proof}
By $\name{Quotient.inductionOn}$ it suffices to treat a class represented by a
bag $b$, where $\name{eval}$ reduces to $\name{PrimeBag.eval}(b)$.  Then
$0 < 1$ and $1 \leq \name{PrimeBag.eval}(b)$ ---
the latter is $\name{one\_le\_eval}$ from \texttt{PrimeTensor/Prime.lean} ---
compose to $0 < \name{PrimeBag.eval}(b)$.
\end{proof}

The lemma is \lean{private}, hence visible only inside this file.  That is
appropriate: it exists to feed the single cancellation step below, and the
public statement of that fact is $\name{one\_le\_eval}$, which is sharper.

\subsection{Cross-product equality is an equivalence}

\begin{lemma}[\lean{same\_refl} and \lean{same\_symm}]\label{MulRat:lem:refl-symm}
$\name{Same}$ is reflexive and symmetric.
\end{lemma}

\begin{proof}
Unfolding the definition from \texttt{PrimeTensor/Ratio.lean},
$\name{Same}(a,a)$ is
$\name{eval}(a.\name{upper}) \cdot \name{eval}(a.\name{lower})
 = \name{eval}(a.\name{upper}) \cdot \name{eval}(a.\name{lower})$,
closed by $\ctor{rfl}$; and $\name{Same}(a,b)$ and $\name{Same}(b,a)$ are the
two sides of one equation of naturals, so symmetry is $\name{Eq.symm}$.  The
two carry \lean{@[refl]} and \lean{@[symm]}.
\end{proof}

\begin{lemma}[Transitivity, \lean{PrimeRatio.same\_trans}]\label{MulRat:lem:trans}
If $\name{Same}(a,b)$ and $\name{Same}(b,c)$ then $\name{Same}(a,c)$.
\end{lemma}

\begin{leancode}
@[trans] theorem same_trans {a b c : PrimeRatio}
    (hab : Same a b) (hbc : Same b c) : Same a c := by
  unfold Same at hab hbc ⊢
  apply Nat.eq_of_mul_eq_mul_right (multiset_eval_pos b.lower)
  calc
    (a.upper.eval * c.lower.eval) * b.lower.eval
        = (a.upper.eval * b.lower.eval) * c.lower.eval := by ac_rfl
    _ = (b.upper.eval * a.lower.eval) * c.lower.eval := by rw [hab]
    _ = (b.upper.eval * c.lower.eval) * a.lower.eval := by ac_rfl
    _ = (c.upper.eval * b.lower.eval) * a.lower.eval := by rw [hbc]
    _ = (c.upper.eval * a.lower.eval) * b.lower.eval := by ac_rfl
\end{leancode}

\begin{proof}
Write $u_x = \name{eval}(x.\name{upper})$ and
$l_x = \name{eval}(x.\name{lower})$.  The hypotheses are
$u_a l_b = u_b l_a$ and $u_b l_c = u_c l_b$, and the goal is
$u_a l_c = u_c l_a$.

The obvious route --- divide the two hypotheses --- is unavailable, since
$\mathbb{N}$ has no division.  Instead one multiplies the goal by $l_b$ and
cancels it afterwards.  By $\name{Nat.eq\_of\_mul\_eq\_mul\_right}$, applied
with the positivity of $l_b$ supplied by Lemma~\ref{MulRat:lem:pos}, it is
enough to prove
\[
  (u_a l_c) \, l_b = (u_c l_a) \, l_b,
\]
and that is the displayed chain: reassociate to expose $u_a l_b$, rewrite it
to $u_b l_a$ by the first hypothesis, reassociate to expose $u_b l_c$, rewrite
it to $u_c l_b$ by the second, and reassociate once more.  Each reassociation
step is closed by $\ctor{ac\_rfl}$, which decides equalities that hold by
associativity and commutativity of multiplication alone.
\end{proof}

\begin{designnote}[Where the missing zero is spent]\label{MulRat:note:cancel}
Cancellation is the whole difficulty, and cancellation on $\mathbb{N}$ is
valid only against a factor that is nonzero.  The factor here is $l_b$, the
evaluation of a prime multiset, and Lemma~\ref{MulRat:lem:pos} makes it
positive by tracing back through $\name{one\_le\_eval}$ to the fact that a
prime bag is a product of atoms each greater than $1$, terminated by the
pivot $\ctor{one}$ rather than by a zero.

Had the carrier admitted a zeroth stage, this step would have needed a side
condition, and the side condition would have propagated into $\name{setoid}$,
into every lift out of the quotient, and into every downstream statement about
$\name{MulRat}$.  Instead there is nothing to discharge: the representation
cannot express the bad case.
\end{designnote}

\begin{definition}[The setoid, \lean{PrimeTensor.PrimeRatio.setoid}]
\label{MulRat:def:setoid}
$\name{setoid} : \name{Setoid}(\name{PrimeRatio})$ is $\name{Same}$ with the
equivalence proof assembled from Lemmas~\ref{MulRat:lem:refl-symm}
and~\ref{MulRat:lem:trans}.
\end{definition}

\subsection{Congruence and the laws at representative level}

\begin{lemma}[\lean{mul\_same} and \lean{inv\_same}]\label{MulRat:lem:congr}
$\name{Same}$ is respected by multiplication and by inversion:
if $\name{Same}(a_1,a_2)$ and $\name{Same}(b_1,b_2)$ then
$\name{Same}(a_1 b_1,\, a_2 b_2)$; and if $\name{Same}(a,b)$ then
$\name{Same}(a^{-1}, b^{-1})$.
\end{lemma}

\begin{proof}
For the first, $\name{upper\_mul}$ and $\name{lower\_mul}$ from
\texttt{PrimeTensor/Ratio.lean} turn the goal into
\[
  (u_{a_1} u_{b_1})(l_{a_2} l_{b_2})
    = (u_{a_2} u_{b_2})(l_{a_1} l_{b_1}),
\]
and regrouping gives
$(u_{a_1} l_{a_2})(u_{b_1} l_{b_2})$, which the two hypotheses rewrite to
$(u_{a_2} l_{a_1})(u_{b_2} l_{b_1})$, which regroups to the right-hand side.
For the second, $\name{inv\_upper}$ and $\name{inv\_lower}$ turn the goal into
$l_a u_b = l_b u_a$, and the hypothesis $u_a l_b = u_b l_a$ gives it after
commuting each side.  Both proofs are chains of $\ctor{ac\_rfl}$ steps around
a single rewrite.
\end{proof}

\begin{lemma}[\lean{one\_upper} and \lean{one\_lower}]\label{MulRat:lem:one-fields}
$(1 : \name{PrimeRatio}).\name{upper} = 1$ and
$(1 : \name{PrimeRatio}).\name{lower} = 1$.
\end{lemma}

\begin{leancode}
@[simp] theorem one_upper : (1 : PrimeRatio).upper = 1 := rfl
@[simp] theorem one_lower : (1 : PrimeRatio).lower = 1 := rfl
\end{leancode}

\begin{proof}
$\ctor{rfl}$: the unit of $\name{PrimeRatio}$ is the literal
$\langle 1, 1 \rangle$, and projecting a literal reduces.
\end{proof}

\begin{lemma}[The laws, up to $\name{Same}$]\label{MulRat:lem:laws-same}
For all $a, b, c : \name{PrimeRatio}$:
\[
  \begin{array}{ll}
    \name{Same}(1 \cdot a,\; a), &
      \text{\lean{one\_mul\_same}} \\
    \name{Same}(a \cdot 1,\; a), &
      \text{\lean{mul\_one\_same}} \\
    \name{Same}((a b) c,\; a (b c)), &
      \text{\lean{mul\_assoc\_same}} \\
    \name{Same}(a b,\; b a), &
      \text{\lean{mul\_comm\_same}} \\
    \name{Same}((a^{-1})^{-1},\; a), &
      \text{\lean{inv\_inv\_same}} \\
    \name{Same}(a \cdot a^{-1},\; 1), &
      \text{\lean{mul\_inv\_same\_one}}
  \end{array}
\]
\end{lemma}

\begin{proof}
Each unfolds $\name{Same}$, rewrites the components with the evaluation lemmas
of \texttt{PrimeTensor/Ratio.lean} --- $\name{upper\_mul}$,
$\name{lower\_mul}$, $\name{inv\_upper}$, $\name{inv\_lower}$ --- together
with Lemma~\ref{MulRat:lem:one-fields} and $\name{PrimeMultiset.eval\_one}$,
and closes the resulting identity of naturals.  The unit laws close by the
unit laws of $\mathbb{N}$; associativity and commutativity close by
$\ctor{ac\_rfl}$; double inversion closes by simplification alone, both sides
becoming literally the same expression once the two swaps cancel; and
$\name{Same}(a a^{-1}, 1)$ reduces to
$u_a l_a = l_a u_a$, which is $\name{Nat.mul\_comm}$.
\end{proof}

Note what these say and do not say.  $\name{PrimeRatio}$ itself does not
satisfy any of the six laws --- $a \cdot 1$ and $a$ are different terms, and
$a \cdot a^{-1}$ is $\langle u_a l_a, l_a u_a \rangle$, not
$\langle 1, 1 \rangle$.  Each law holds only after $\name{Same}$ has been
applied, which is exactly the input the quotient needs.

\subsection{The quotient}

\begin{definition}[\lean{PrimeTensor.MulRat}]\label{MulRat:def:mulrat}
$\name{MulRat} \defeq \name{Quotient}(\name{PrimeRatio.setoid})$, an
\lean{abbrev} and hence reducible.  $\name{ofRatio}(q)$ is the class of $q$,
and $\name{one} \defeq \name{ofRatio}(1)$.
\end{definition}

\begin{definition}[Induced operations]\label{MulRat:def:ops}
The declarations \lean{MulRat.mul} and \lean{MulRat.inv} lift multiplication
and inversion through the quotient; \lean{One}, \lean{Mul} and \lean{Inv}
instances give them their notation.
\end{definition}

\begin{leancode}
def mul (a b : MulRat) : MulRat :=
  Quotient.liftOn₂ a b
    (fun x y => ofRatio (x * y))
    (by
      intro a₁ b₁ a₂ b₂ ha hb
      change PrimeRatio.Same a₁ a₂ at ha
      change PrimeRatio.Same b₁ b₂ at hb
      change Quotient.mk PrimeRatio.setoid (a₁ * b₁) =
        Quotient.mk PrimeRatio.setoid (a₂ * b₂)
      apply Quotient.sound
      exact PrimeRatio.mul_same ha hb)
\end{leancode}

Each lift carries a well-definedness obligation, and each is exactly one of
the congruence lemmas: $\name{mul\_same}$ for the binary lift,
$\name{inv\_same}$ for the unary one, in both cases followed by
$\name{Quotient.sound}$ to convert the resulting $\name{Same}$ into an
equality of classes.

\begin{theorem}[$\name{MulRat}$ is a commutative group]\label{MulRat:thm:group}
For all $a, b, c : \name{MulRat}$,
\[
  1 \cdot a = a, \quad
  a \cdot 1 = a, \quad
  (ab)c = a(bc), \quad
  ab = ba, \quad
  (a^{-1})^{-1} = a, \quad
  a \cdot a^{-1} = 1,
\]
these being the declarations \lean{one\_mul}, \lean{mul\_one},
\lean{mul\_assoc}, \lean{mul\_comm}, \lean{inv\_inv} and \lean{mul\_inv}.
\end{theorem}

\begin{proof}
Every one has the same three-step shape.  By $\name{Quotient.inductionOn}$
(or its two- and three-variable forms) it suffices to prove the statement for
classes of representatives; there the goal becomes an equality of two classes;
$\name{Quotient.sound}$ reduces that to the corresponding $\name{Same}$
statement, which is the matching item of Lemma~\ref{MulRat:lem:laws-same}.
\end{proof}

\begin{remark}[What is registered, and what is not]\label{MulRat:rem:instances}
The four laws that make $\name{MulRat}$ a group are theorems, but no
\lean{Group} or \lean{CommGroup} instance is declared --- only \lean{One},
\lean{Mul} and \lean{Inv}.  So Mathlib's group lemmas are not available for
$\name{MulRat}$ downstream, and each fact must be taken from this file by
name.  Note also that $a^{-1} \cdot a = 1$ is not stated, although it follows
from $\name{mul\_inv}$ and $\name{mul\_comm}$; and that
$\name{one\_mul}$, $\name{mul\_one}$, $\name{inv\_inv}$ and $\name{mul\_inv}$
carry \lean{@[simp]} while $\name{mul\_assoc}$ and $\name{mul\_comm}$ do not,
as is usual for laws that would loop the simplifier.
\end{remark}

\begin{remark}[What the quotient is]\label{MulRat:rem:iso}
Sending $a$ to
$\name{eval}(a.\name{upper}) / \name{eval}(a.\name{lower})$ maps
$\name{PrimeRatio}$ onto $\mathbb{Q}_{>0}$: it is surjective because every
positive rational is a quotient of two positive integers and every positive
integer is a product of primes, and by the cross-product characterisation in
\texttt{PrimeTensor/Ratio.lean} two ratios have the same image precisely when
$\name{Same}$ relates them.  So $\name{MulRat}$ is isomorphic to
$(\mathbb{Q}_{>0}, \cdot, 1)$ as a group, and
Theorem~\ref{MulRat:thm:group} is the barcode-level reflection of that.  The
isomorphism is stated here for the reader; the file constructs no such map,
and consistently with its header no value of type $\mathbb{Q}$ or $\mathbb{R}$
appears in it.
\end{remark}

\begin{remark}[The order does not descend here]\label{MulRat:rem:no-order}
\texttt{PrimeTensor/Ratio.lean} also defined $\name{Lt}$ and $\name{Within}$.
This file proves nothing about either, and $\name{MulRat}$ carries no order:
there is no \lean{LT MulRat} instance and no lemma showing $\name{Lt}$
respects $\name{Same}$, though it does.  Only the multiplicative structure is
transported across the quotient.
\end{remark}

\subsection*{Summary of the correspondence}

\begin{center}
\small
\begin{tabular}{@{}lll@{}}
\hline
Lean declaration & Kind & Here \\
\hline
\lean{PrimeRatio.multiset\_eval\_pos} & private thm & Lem.~\ref{MulRat:lem:pos} \\
\lean{PrimeRatio.same\_refl}      & refl thm   & Lem.~\ref{MulRat:lem:refl-symm} \\
\lean{PrimeRatio.same\_symm}      & symm thm   & Lem.~\ref{MulRat:lem:refl-symm} \\
\lean{PrimeRatio.same\_trans}     & trans thm  & Lem.~\ref{MulRat:lem:trans} \\
\lean{PrimeRatio.setoid}          & definition & Def.~\ref{MulRat:def:setoid} \\
\lean{PrimeRatio.mul\_same}       & theorem    & Lem.~\ref{MulRat:lem:congr} \\
\lean{PrimeRatio.inv\_same}       & theorem    & Lem.~\ref{MulRat:lem:congr} \\
\lean{PrimeRatio.one\_upper}      & simp thm   & Lem.~\ref{MulRat:lem:one-fields} \\
\lean{PrimeRatio.one\_lower}      & simp thm   & Lem.~\ref{MulRat:lem:one-fields} \\
\lean{PrimeRatio.one\_mul\_same}  & theorem    & Lem.~\ref{MulRat:lem:laws-same} \\
\lean{PrimeRatio.mul\_one\_same}  & theorem    & Lem.~\ref{MulRat:lem:laws-same} \\
\lean{PrimeRatio.mul\_assoc\_same}& theorem    & Lem.~\ref{MulRat:lem:laws-same} \\
\lean{PrimeRatio.mul\_comm\_same} & theorem    & Lem.~\ref{MulRat:lem:laws-same} \\
\lean{PrimeRatio.inv\_inv\_same}  & theorem    & Lem.~\ref{MulRat:lem:laws-same} \\
\lean{PrimeRatio.mul\_inv\_same\_one} & theorem & Lem.~\ref{MulRat:lem:laws-same} \\
\lean{PrimeTensor.MulRat}         & abbrev     & Def.~\ref{MulRat:def:mulrat} \\
\lean{MulRat.ofRatio}             & definition & Def.~\ref{MulRat:def:mulrat} \\
\lean{MulRat.one}                 & definition & Def.~\ref{MulRat:def:mulrat} \\
\lean{MulRat.mul}                 & definition & Def.~\ref{MulRat:def:ops} \\
\lean{MulRat.inv}                 & definition & Def.~\ref{MulRat:def:ops} \\
\lean{One}, \lean{Mul}, \lean{Inv} & instances & Def.~\ref{MulRat:def:ops} \\
\lean{MulRat.one\_mul}            & simp thm   & Thm.~\ref{MulRat:thm:group} \\
\lean{MulRat.mul\_one}            & simp thm   & Thm.~\ref{MulRat:thm:group} \\
\lean{MulRat.mul\_assoc}          & theorem    & Thm.~\ref{MulRat:thm:group} \\
\lean{MulRat.mul\_comm}           & theorem    & Thm.~\ref{MulRat:thm:group} \\
\lean{MulRat.inv\_inv}            & simp thm   & Thm.~\ref{MulRat:thm:group} \\
\lean{MulRat.mul\_inv}            & simp thm   & Thm.~\ref{MulRat:thm:group} \\
\hline
\end{tabular}
\end{center}

\noindent
Every declaration of \texttt{PrimeTensor/MulRat.lean} appears above.  The file
contains no \lean{sorry}, no axiom, and no unproved named proposition.
Remarks~\ref{MulRat:rem:iso} and~\ref{MulRat:rem:no-order} are stated for the
reader and are not declarations of the file.
